\documentclass[anon]{alt2027} % submission template (anonymized)
%\documentclass[final]{alt2027} % camera-ready template (de-anonymized)

% The following packages will be automatically loaded:
% amsmath, amssymb, natbib, graphicx, url, xcolor, hyperref, algorithm2e
% Keep alt2027.cls, jmlr.cls, and jmlrutils.sty in the same directory as
% this file. Compile with latexmk -pdf alt2027-sample.tex or your usual
% LaTeX workflow.


% caligraphic letters
  \newcommand{\Acal}{{\mathcal A}}
  \newcommand{\A}{{\mathcal A}}
  \newcommand{\Bcal}{{\mathcal B}}
  \newcommand{\B}{{\mathcal B}}
  \newcommand{\Ccal}{{\mathcal C}}
  \newcommand{\C}{{\Ccal}}
  \newcommand{\Dcal}{{\mathcal D}}
  \newcommand{\Fcal}{{\mathcal F}}
  \newcommand{\F}{{\mathcal F}}
  \newcommand{\Hcal}{{\mathcal H}}
  \renewcommand{\H}{{\Hcal}}
  \newcommand{\Ical}{{\mathcal I}}
  \newcommand{\Pcal}{{\mathcal P}}
  \renewcommand{\P}{{\mathcal P}}
  \newcommand{\Qcal}{{\mathcal Q}}
  \newcommand{\Ocal}{{\mathcal O}}
  \newcommand{\Ncal}{{\mathcal N}}
  \newcommand{\Tcal}{{\mathcal T}}
  \newcommand{\T}{{\mathcal T}}
  \newcommand{\Ucal}{{\mathcal U}}
  \newcommand{\Vcal}{{\mathcal V}}
  \newcommand{\V}{{\mathcal V}}
  \newcommand{\Scal}{{\mathcal S}}
  \newcommand{\Wcal}{{\mathcal W}}
  \newcommand{\Xcal}{{\mathcal X}}
  \newcommand{\X}{{\mathcal X}}
  \newcommand{\Z}{{\mathcal Z}}
  \newcommand{\Y}{{\mathcal Y}}
  \newcommand{\Ycal}{{\mathcal Y}}
  \newcommand{\Zcal}{{\mathcal Z}}
  \newcommand{\ol}{\overline}
  \newcommand{\zero}{\bf{0}}

  \newcommand{\Lo}{{\mathcal L}}
  \renewcommand{\Re}{{\mathcal R}}
  \newcommand{\Mi}{{\mathcal M}}


 \newcommand{\ie}{\textsl{i.e.\ }} % id est
 \newcommand{\iid}{\textsl{i.i.d.\ }} % id est
 \newcommand{\eg}{e.g.\ }          % e.g.
 \newcommand{\ex}{\mathbb{E}}
 \newcommand{\reals}{\mathbb{R}}
 \newcommand{\naturals}{\mathbb{N}}

 %certificates and debaters
  \newcommand{\dt}{\texttt{Proposer}\xspace}
  \newcommand{\df}{\texttt{Responder}\xspace}
  \renewcommand{\c}{c}
  \newcommand{\ct}{{p}}
  \newcommand{\cf}{{r}}
  \newcommand{\hy}{{\hat{y}}}
  \newcommand{\ty}{{y^*}}


% comments
\newcommand{\todo}[1]{{\color{red}todo: #1}}
\newcommand{\ruth}[1]{{\color{purple}Ruth: #1}}
\newcommand{\tosca}[1]{{\color{teal}Tosca: #1}}
\newcommand{\sivan}[1]{{\color{orange}todo: #1}}
\newcommand{\vinayak}[1]{{\color{blue}Vinayak: #1}}


\title[Learnability via Debate]{Towards a Theory of Learnability via Debate}
\usepackage[T1]{fontenc}
\usepackage{times}
% Use \Name{Author Name} to specify the name.
% If the surname contains spaces, enclose the surname
% in braces, e.g. \Name{John {Smith Jones}} similarly
% if the name has a "von" part, e.g \Name{Jane {de Winter}}.
% If the first letter in the forenames is a diacritic
% enclose the diacritic in braces, e.g. \Name{{\'E}louise Smith}

% Two authors with the same address
% \altauthor{\Name{Author Name1} \Email{abc@sample.com}\and
%  \Name{Author Name2} \Email{xyz@sample.com}\\
%  \addr Address}

% Three or more authors with the same address:
% \altauthor{\Name{Author Name1} \Email{an1@sample.com}\\
%  \Name{Author Name2} \Email{an2@sample.com}\\
%  \Name{Author Name3} \Email{an3@sample.com}\\
%  \addr Address}

% Authors with different addresses:
\altauthor{%
 \Name{Author Name1} \Email{abc@sample.com}\\
 \addr Address 1
 \AND
 \Name{Author Name2} \Email{xyz@sample.com}\\
 \addr Address 2%
}

\begin{document}

\maketitle

\begin{abstract}%
  We initiate the formal study of learning via debate.%
\end{abstract}

\begin{keywords}%
  Learning via debate, online learning, Littlestone tree
\end{keywords}

\section{Introduction}
As AI models become smarter, how do we ensure that they continue to act in the interest of humans? A long line of work~\citep{} has argued that this is a hard problem, and recent events such as the infamous HuggingFace incident~\citep{} have demonstrated the difficulties in practice. 

One way to study this problem more formally has been termed \emph{scalable oversight}. Here, a smarter-than-human AI is typically modelled as an algorithm that can solve computational problems in some large complexity class, and a human is modelled as an algorithm in a much smaller complexity class. It is assumed that given a context $x$, computing a human-aligned action $y$ is hard for the human but possible for the AI. How can the human verify whether the AI is, in fact, calculating $y$, as opposed to some other bad action $y'$? Of course, without any extra assumptions this is impossible. The scalable oversight literature assumes that the AI sits within some level of the polynomial hierarchy, thus providing a means to \emph{certify} whether AI's solution is correct. For example, for the first level of the hierarchy, one assumes the AI to be in $\operatorname{NP} \cap \operatorname{coNP}$ and the human in P, thus implying that there exists a certificate $c$ that the human can look at to verify whether $y$ is the correct answer to $x$. For the second level, this implies there exists a certificate $c_1$ such that for all certificates $c_2$, the human's response to seeing $x, y, c_1$, and $c_2$ is to judge $y$ to be the correct answer to $x$. For the $k^\text{th}$ level, this looks like debate where two debaters alternately produce certificates in a way that the first one has a winning strategy if and only if $y$ is the correct answer. 


\section{Formal Setup}
% We let $\X$ denote a \emph{domain}, $\Y = \{0,1\}$ a binary \emph{label space}, and $\C = \{0,1\}$ a binary \emph{certificate space}.  
We let $\X$ denote a \emph{domain}, $\Y $ a \emph{label space}, and $\C$ a \emph{certificate space}.  
A \emph{hypothesis} is a function $h:\X\to\Y$, and we use $\H$ to denote a \emph{hypothesis class}.
A \emph{debate of length} $n$ is a sequence $(\c_1, \c_2, \ldots, \c_{2n})$ of $2n$ certificates $c_i\in\C$. We think of these certificates as being produced alternately by two debaters. That is, we think of the certificates $c_i$ of odd index as being produced by the first debater, which we will call the \dt, and the certificates $c_i$ of even index by the second debater, which we will refer to as the \df. To clarify the role of the arguments in the debate, we usually will use the notation $(\ct^1, \cf^1, \ct^2, \cf^2, \ldots,  \ct^n, \cf^n)$ for a debate of length $n$

The debate is moderated through a \emph{verifier} $v: \X\times \C^{2n} \to \Y$ which takes in a domain point and a debate of length $n$ and produces a label for $x$ given the arguments in the debate. The following definition specifies the relationship between hypotheses and verifiers in our framework.

\begin{definition}
Let $h:\X\to\Y$ be a hypothesis. We say that verifier $v: \X\times \C^{2n} \to \Y$ \emph{validates} hypothesis $h$ through debate of length $n$ if 
% \begin{align*}
% ~\forall~ x\in \X ~\exists~ \ct^1 \in \C ~\forall~ \cf^1 \in \C   ~\exists~ \ct^2 \in \C ~\forall~ & \cf^2 \in \C \ldots ~\exists~ \ct^n \in \C ~\forall~ \cf^n ~:~ \\
% & v(x, \ct^1, \cf^1, \ct^2, \cf^2, \ldots,  \ct^n, \cf^n) = h(x)
% \end{align*}
\begin{align*}
~\forall~ x\in \X ~\exists~ \ct^1 \in \C ~\forall~ \cf^1 \in \C &  ~\exists~ \ct^2 \in \C ~\forall~  \cf^2 \in \C \ldots ~\exists~ \ct^n \in \C ~\forall~ \cf^n \in \C ~:~ \\
& v(x, \ct^1, \cf^1, \ct^2, \cf^2, \ldots,  \ct^n, \cf^n) = y \iff h(x) = y
\end{align*}
% \vinayak{
% This ensures that the situation where there is a valid debate transcript for both labels does not happen.
% }
In this case, we also say that $v$ is a verifier for hypothesis $h$. We call a verifier \emph{valid} if there exists a hypothesis $h$ that is validated by $v$.
\end{definition}
Note that a given valid verifier $v$ validates only one unique hypothesis $h$ in the above sense. On the other hand, a given hypothesis $h$ may be validated by multiple verifiers $v$.\ruth{This is still true? Does it matter? Or should we assume $\V$ is pruned in a way that there is a one to one correspondence between $\H$ and $\V$?}
% ~\vinayak{Oh I see. I have been thinking of verifiers as uniquely defining a hypothesis. Given an $x$ and a $v$, I think only one unique label should meet the condition above, and that label will define a unique hypothesis. I think Tosca is saying the same thing below.}. 
We will call a set of valid verifiers a \emph{verifiers class}, and use notation $\V$ to denote a verifier class. Further, we let $\H_\V$ denote the hypothesis class of all functions that are validated by verifiers $v\in\V$. If clear from context, we may simply write $\H$ for $\H_\V$.
% \tosca{I think one thing that is still missing here is that we want the label that is certified to be unique. In the current definition there exists a hypothesis that verfies \emph{all hypotheses}. I suggest we add "soundness" of a verifier somewhere (i.e. it for every $x$ there exists no more than one label $y$ that $v$ verfifies)}

\paragraph{Learning via debate} The goal of \emph{learning via debate} is to recover an underlying \emph{true hypothesis}, which we denote by $h^*$. We assume that the two debaters interact with a verifier $v^*\in\V$. In the \emph{realizable case} the verifier $v^*$ validates the true hypothesis $h^*$ (and thus $h^*\in\H = \H_\V$). Of course, debaters \dt and \df don't know the identity of $h^*$ nor the identity of $v^*\in\V$. They only have access to the verifier class $\V$, and thereby also to its induced hypothesis class $\H_\V$.

\paragraph{The debate protocol}

\begin{itemize}
\item For rounds $t = 1\ldots T$ repeat:
\begin{enumerate}
 \item The environment poses an $x_t\in\X$
 \item Debater \dt hypothesizes label $\hy_t$
 \item For $i = 1 \ldots n$ repeat:
 \begin{itemize}
  \item Debater \dt brings argument $\ct^i_t$ (in support of label $\hy_t$)
  \item Debater \df brings argument $\cf^i_t$ (aiming to falsify label $\hy_t$)
 \end{itemize}
 \item The verifier $v^*$ reveals $y_t = v(x, \ct^1_t, \cf^1_t, \ct^2_t, \cf^2_t, \ldots,  \ct^n_t, \cf^n_t)$\\
       The environment records (but does not reveal) $y^*_t = h^*(x_t)$
\end{enumerate}
\item The debate is evaluated through the \emph{number of mistakes} and the \emph{debate regret}
\begin{align*}
 \Mi(T) = \sum_{t = 1}^T [y_t \neq y^*_t]\quad\text{define regret in comparison to optimal debaters(?)}
\end{align*}
\item Debater \dt is evaluated through its \emph{loss} and its \emph{regret}
\begin{align*}
 \Lo_T(T) = \sum_{t = 1}^T [y_t \neq \hy_t]\quad\text{and}\quad \Re_T(T) = \Lo_T(T) - \#\text{unachievable rounds for \dt}
\end{align*}
\ruth{In the realizable case all rounds are achievable for \dt, but probably not in the agnostic case?} \tosca{We got into a debate about how to formalize this in our meeting in Hamilton last week as well. In other settings we often just deducted $k$ (the number of exceptions), which may be an overestimation of non-achievable rounds. }
\item Debater \df is evaluated through \emph{loss} and its \emph{regret}
\begin{align*}
 \Lo_F(T) = \sum_{t = 1}^T [y_t = \hy_t]\quad\text{and}\quad \Re_F(T) = \Lo_F(T) - \#\text{unachievable rounds for \df}
\end{align*}
\end{itemize}
\ruth{We might want to make the presentation of the protocol more compact; still need to define unachievable rounds properly}
\vinayak{Things are a bit more complicated in the agnostic case because two different kinds of things make it impossible to achieve 0 loss: 1) if there is misspecification, i.e., no hypothesis in the class assigns the true label, and b) if its an anachievable round, i.e., \dt has picked the ``correct'' certificate and so \df cannot possibly win.}

The number of rounds may be finite ($T\in\naturals$) or infinite ($T = \infty$).
We note that the evaluation of the two debaters is distinct from the evaluation of the overall debate. Debater \dt is evaluated against its aim to defend its hypothesized label $\hy_t$ in each round $t$. Debater \df on the other hand is evaluated against its aim to falsify the hypothesis that the first debater proposed. That is, debater \df, is aiming to set its certificates in such a way that the verifier will conclude with a label $y_t$ that is not the hypothesized label $\hy_t$. The value of the overall debate that the two debaters produce is measured by assessing whether the sequence of produced arguments/certificates actually lead the verifier to conclude the correct label $y^*_t$ in each round.

We will study for which verifier classes $\V$ the above protocol leads to a finite mistake bound (or sublinear regret) for the debate evaluation. We will distinguish between the \emph{cooperative case} and the \emph{combative case}. In the cooperative case both debaters are acting cooperatively so as to minimize the debate evaluation measures. We can think of this as a scenario where debater \df provides its certificates in good faith ``for the sake of the argument'' so as to lead the verifier to the correct conclusions having taking into account relevant potential counter-arguments. In contrast, in the combative case, both debaters are acting in a way to minimize their own evaluation measures respectively. Thus in this case, even debater true may put up hypothesized labels $\hy_t$ that it knows it can argue for even through it is already aware the this does not correspond to the underlying truth $y^*_t$. We study whether the true underlying hypotheses can be recovered in both scenarios.

\vinayak{I'm not sure that I like the names ``cooperative'' and ``combative''. The point of studying the cooperative version was to understand the best way to use access to this kind of verifiers. The ``debate'' thesis is that the best way to use such a verifier is to set up two debaters and let them argue against each other. But we don't think that should be the case. So the question is, if we don't restrict ourselves to set up these debaters and simply understand the best way to use these verifiers, what's the best that can be done? So ``cooperative'' is really just more general and ``combative'' is more restrictive. This should be reflected in the names somehow.}

\section{The Cooperative Case}
In this section, we start by exploring the cooperative case of learning via debate. To analyze which classes $\V$ of verifiers allow for learning with finite mistake bound, we will define a tree structure that represents all potential runs of the debate protocol for a given starting instance $x_0$. 

\begin{definition}
Given a debate length $n$, we define a \emph{length-$n$-debate tree} to be the directed tree $\T$ with the following properties:
\begin{enumerate}
    \item Every node $\alpha$ in $\T$ is a pair of a domain point and a label $\alpha=(x_\alpha, y_\alpha)$.
    \item The root node is denoted by $\rho$. It's domain point is $x_\rho$ and it's label is $y_\rho = \sqcup$, where $\sqcup\notin\Y$. \item The root node $\rho$ is the only node in $\T$ with associated label $y_\rho \notin\Y$. For all nodes $\alpha\neq \rho$, we have $y_\alpha\in \Y$.
    \item Edges $(\alpha, \beta)$ in $\T$ are directed away from the root node, and annotated with a label $\ell^\alpha_\beta = (\hy^\alpha_\beta, \bar{c}^\alpha_\beta)$, where $\bar{c} = (\ct^1_{\alpha\beta}, \cf^1_{\alpha\beta}, \ct^2_{\alpha\beta}, \cf^2_{\alpha\beta}, \ldots,  \ct^n_{\alpha\beta}, \cf^n_{\alpha\beta})$ is a debate of length $n$. 
    % We will denote edges by triples
    % \[
    % (\alpha, \ell, \beta) = ((x_\alpha, y_\alpha), (\hy^\alpha_\beta, \bar{c}^\alpha_\beta),  (x_\beta, y_\beta)).
    % \]
    \item For each node $\alpha$ and each possible edge annotation $(\hy, \bar{c})$, there is exactly one child node $\beta$ to $\alpha$ with $\ell^\alpha_\beta = (\hy, \bar{c})$
\end{enumerate}
We let $\Delta_n$ denote the set of all length-$n$-debate trees $\T$.
\end{definition}



\section{The Combative Case}
% Authors are strongly encouraged to include an anonymous AI Disclosure section
% if generative AI was used substantively. Uncomment and replace the example:
% \section*{AI Disclosure}
% We used [tool] to [briefly explain how it was used].

% Acknowledgments---Will not appear in anonymized version
\acks{We thank our colleagues and funding agencies;
  \texttt{\textbackslash documentclass[anon]\{alt2027\}}
automatically hides this text.}

% Replace yourbibfile with the name of your BibTeX database and uncomment:
% \bibliography{yourbibfile}

\appendix

% \crefalias{section}{appendix} % uncomment if you are using cleveref

\section{Deferred proofs}

Here is the proof of the main lemma from the paper body.

\begin{proofof}{Theorem 1}
  Follows from the concavity of $\ln(\cdot)$ and Jensen's inequality.
\end{proofof}

\end{document}
